\documentclass[11pt]{article}
\usepackage{kept-report}
\title{Fairness report}
\subtitle{Does Kept treat protected groups differently? Checked on tonight's 96,000 cards}
\reportnumber{Governance · fairness}
\runninghead{Fairness}
\date{10 October 2026}

\begin{document}
\makereporttitle

\statrow{\stat{0.87–1.00}{any-contact ratio, every attribute (flag below 0.80)}\stat{0 of 96,000}{cards change when every protected attribute is scrambled}\stat{11}{attributes checked, from the data contract}\stat{7}{flags, all explained below}}

\takeaway{\textbf{The short version.} For every protected attribute, the group that gets contacted least is within
0.87–1.00 of the group that gets contacted most, so nothing comes near the four-fifths line. The 7 flags sit on two outcomes.
\textbf{Calls} are rare (2.9\% of cards), so a handful of calls in a small group flips the ratio. \textbf{Hardship review}
lands more on customers who carry the vulnerability flag, which is the protection doing its job. One gap needs a human look:
separated customers get a hardship review 51\% of the time against 39–40\% for everyone else. Most of that is the
vulnerability flag, and §4 shows the rest.}

\section{How we checked}

\begin{center}
\begin{tikzpicture}[node distance=0.45cm, every node/.style={font=\small\sffamily, align=center},
  box/.style={draw=rule, fill=soft, rounded corners=3pt, inner sep=6pt, text width=2.9cm},
  hot/.style={box, fill=kept!7, draw=kept!40},
  go/.style={-{Stealth[length=2mm]}, draw=muted, line width=0.7pt}]
\node[box] (c) {\textbf{Tonight's cards}\\96,000 cases, one action each};
\node[hot, right=of c] (r) {\textbf{c360\_restricted}\\compliance-only; the model can't read it};
\node[box, right=of r] (j) {\textbf{Join on golden ID}\\one row per card, plus its attributes};
\node[box, right=of j] (m) {\textbf{Rate per group}\\lowest ÷ highest, groups of 50+};
\draw[go] (c) -- (r); \draw[go] (r) -- (j); \draw[go] (j) -- (m);
\end{tikzpicture}
\end{center}

The attributes are the data contract's protected and proxy list (gender, citizenship, marital status, household,
dependants, accessibility, vulnerability type, benefit deposits, accent region, newcomer programme). Each is checked on
three outcomes: \emph{any contact action}, \emph{call} and \emph{hardship review}. The test is the four-fifths rule: the
lowest group's rate divided by the highest, flagged below 0.80. Groups under 50 people are left out, because their rates are
noise.

Keeping these attributes out of decisions is enforced in three places, not promised:
\begin{itemize}
\item \textbf{Storage.} They live only in \code{gold.c360\_restricted}, which only the compliance role can read.
\item \textbf{CI.} A lineage test fails the build if any model feature traces back to one of them, or to a proxy such as postal
code, age or language.
\item \textbf{Tests.} A curated suite checks the trained model's actual input columns, and checks that the Ask refuses
questions that rank or filter people by them.
\end{itemize}

\section{The overview}

\fig[0.95\linewidth]{f_ratios.png}{Each dot is the lowest group's rate ÷ the highest group's rate, for one attribute and one outcome.}

\textbf{What the chart says.} The blue dots (any contact action) form a tight column at 0.87–1.00. Whichever group
you're in, you're about equally likely to get a next step tonight. The green dots (hardship review) are almost all
above 0.80 too. The orange dots (calls) are the spread-out ones, and the flags sit there. The next section explains why
that spread is noise rather than bias.

\begin{table}[H]\small\centering
\begin{tabular}{@{}lrrrr@{}}
\toprule
\textbf{Attribute} & \textbf{Any contact} & \textbf{Call} & \textbf{Hardship review} & \textbf{People known} \\
\midrule
\code{customer\_accent\_region} & 0.87 & \textcolor{kept}{0.57} & 0.83 & 1,214 \\
\code{vulnerability\_type} & 0.91 & \textcolor{kept}{0.36} & 0.91 & 13,696 \\
\code{dependants\_count} & 0.94 & 0.86 & 0.81 & 65,319 \\
\code{accessibility\_type} & 0.94 & \textcolor{kept}{0.51} & 0.82 & 2,606 \\
\code{marital\_status} & 0.96 & \textcolor{kept}{0.73} & \textcolor{kept}{0.76} & 85,315 \\
\code{household\_size} & 0.97 & \textcolor{kept}{0.78} & 0.93 & 65,259 \\
\code{gender\_code} & 0.98 & \textcolor{kept}{0.72} & 0.99 & 91,611 \\
\code{citizenship\_status} & 0.99 & 0.85 & 0.98 & 91,609 \\
\code{accessibility\_needs\_flag} & 0.99 & 0.93 & 0.97 & 92,542 \\
\code{ccb\_deposit\_flag} & 1.00 & 0.95 & 0.98 & 44,864 \\
\code{newcomer\_program\_flag} & 1.00 & 0.92 & 0.98 & 92,542 \\
\bottomrule
\end{tabular}
\caption{Ratios per attribute; red = below 0.80. A flag is a question for compliance, not a verdict.}
\end{table}

\section{Why the call flags are small numbers}

Only 2,824 of 96,000 cards are calls. Kept sends cheap digital nudges first, and calls go to cases where a human
conversation is worth it. With a base rate that low, a ratio moves a lot on very few people:

\begin{itemize}
\item \textbf{Gender (0.72).} The X group (936 people) got 21 calls, or 2.2\%. The U group (895) got 28, or 3.1\%. F and M sit at 3.1\% and
3.0\%. About 3 more calls to X would close the flag.
\item \textbf{Vulnerability type (0.36).} The lowest group, cognitive decline, has 166 people and 1 call. Vulnerable customers
get a hardship review \emph{instead of} a call (next section), so their call rate is near zero across the board.
\item \textbf{Accent region (0.57).} It's only known for 1,214 people, the ones with a recorded call. The groups are 67–502 people, so
the rates swing on a few calls.
\end{itemize}

\insight{What we'd do about it}{Nothing to the model. The four-fifths rule is meant for outcomes where both groups have
enough events. We'd add a minimum-events rule to the nightly check (for example, at least 20 calls in the smaller group)
so compliance sees flags that mean something.}

\section{Group by group}

\fig[0.98\linewidth]{f_groups.png}{Share of cases with any contact action (blue) and with a hardship review (green), largest groups first.}

\textbf{What the chart says.} Within each attribute the blue bars are all the same length, at about 78–80\%. Gender, citizenship,
household and accessibility show no visible difference on either outcome. Two panels stand out:

\begin{itemize}
\item \textbf{Vulnerability type.} The green bars sit near 70\%, against about 40\% everywhere else. Every customer with a vulnerability
type also carries the vulnerability \emph{flag}, a safeguard in the data dictionary rather than a protected attribute, and the
flag routes a case to a protective conversation before anything is sent. Of the 14,232 flagged customers, 100\% are treated as
hardship and 68\% (9,661) get a hardship review tonight. Of the rest, 2,176 get ``wait, no contact beats waiting today'', 207 get a call,
and the other 2,188 are on insolvency or cease holds or otherwise get no contact.
\item \textbf{Marital status.} Separated customers get a hardship review 51\% of the time, against 39–40\% for every other status.
We checked why. \textbf{30\%} of separated customers carry the vulnerability flag, against 12–15\% of the other statuses. Take
the flagged customers out and the gap shrinks to 42.5\% against 34–35\%, a ratio of about 0.80. What's left lines up with
hardship signals from their notes and cash flow. Separation is often a money shock, so that's plausible, and the model never sees
marital status. We still list it in §6 for compliance to review.
\end{itemize}

\section{The invariance test: scramble them and rescore}

Ratios show outcomes are even; they can't prove the attributes are \emph{unused}. So we tested that directly, the way
an auditor would: change the attributes and see whether any decision moves.

\begin{center}
\begin{tikzpicture}[node distance=0.4cm, every node/.style={font=\small\sffamily, align=center},
  box/.style={draw=rule, fill=soft, rounded corners=3pt, inner sep=6pt, text width=2.6cm},
  hot/.style={box, fill=kept!7, draw=kept!40},
  go/.style={-{Stealth[length=2mm]}, draw=muted, line width=0.7pt}]
\node[box] (w) {\textbf{Two copies}\\of tonight's warehouse};
\node[hot, right=of w] (s) {\textbf{Scramble}\\shuffle each attribute across all rows, in silver};
\node[box, right=of s] (b) {\textbf{Rebuild}\\gold, hardship, features (dbt)};
\node[box, right=of b] (r) {\textbf{Rescore}\\all 96,000 cases, same clock};
\node[hot, right=of r] (c) {\textbf{Compare}\\every card, field by field};
\draw[go] (w) -- (s); \draw[go] (s) -- (b); \draw[go] (b) -- (r); \draw[go] (r) -- (c);
\end{tikzpicture}
\end{center}

\begin{table}[H]\small\centering
\begin{tabularx}{\linewidth}{@{}>{\raggedright\arraybackslash}p{3.6cm}rrL@{}}
\toprule
\textbf{Scrambled} & \textbf{Columns} & \textbf{Cards changed} & \textbf{What changed} \\
\midrule
The brief's protected attributes & 24 & \textbf{0 of 96,000} & nothing: action, tier, reasons, probabilities, hardship, hold-out all identical \\
\quad plus every proxy without an operational use & 43 & 299 & all through one rule: 282 cases \emph{gained} hardship protection, 0 lost it; 17 already in hardship gained a cash-flow line on the banner \\
\bottomrule
\end{tabularx}
\caption{Counterfactual invariance on tonight's book. The scramble reached gold: 50\% of gender codes, 75\% of marital statuses
and 72\% of household sizes in \code{c360\_restricted} changed.}
\end{table}

\paragraph{The one proxy that moves something} All 299 changes trace to \code{ei\_deposit\_flag}. The cash-flow hardship rule
reads it: Employment Insurance deposits that started in the last 60 days mark hardship \emph{possible}. POL-COLL-004 v4.2 names
``waiting for Employment Insurance'' as financial hardship, and the brief says to flag it. The flag can only make a case
gentler. A hardship review replaces a call or a payment link, and it never takes protection away. EI stays on the contract's
proxy list, so no model feature can ever read it (the CI lineage test fails the build if one does). The contract now records this
single use under \code{protectiveUses}, alongside \code{vulnerability\_flag}.

\paragraph{What we left alone, and why} Contact language, address and postal code (time zone, so calling hours) and date of birth
(identity matching, done before gold) are proxies the bank needs in order to contact the right person correctly. They are listed
in the script and kept out of the model's features by the same lineage test.

\insight{A bug the test found}{The first run moved 4 probabilities (no actions or tiers) with \emph{nothing} protected involved:
two trend features are regression slopes whose float sums change in the last bits with row order. A slope of
$2\times10^{-15}$ against $-4\times10^{-16}$ falls either side of a tree split. We now round them to 9 decimals, so a rebuild is
bit-for-bit identical, and the protected run then came back at exactly zero.}

\section{What we read from it, and what's next}

\begin{enumerate}[leftmargin=1.5em]
\item \textbf{Contact is even.} Every attribute sits at 0.87 or above for \emph{any contact action}, and most at 0.93 or above.
\item \textbf{Call flags are small numbers, not bias.} Add a minimum-events rule so the nightly check stops raising them.
\item \textbf{More hardship reviews for vulnerable customers is by design.} The vulnerability flag protects; it never scores.
\item \textbf{Review: separated customers.} After the vulnerability flag, a gap of about 8 points remains, which comes from hardship
signals. We'd ask compliance to sample 50 of these cards and confirm the signals are real.
\item \textbf{Blank attributes.} Customers with no marital status on file (10,685) get fewer actions of any kind (54\% vs 79\%) and fewer hardship reviews (27\% vs 40\%).
That's a data-coverage question for the data team, not a fairness fix.
\end{enumerate}

\reproduce{\path{scripts/invariance_test.py} (protected, then all) scrambles, rebuilds and rescores two warehouse copies and
writes \path{data/exports/reports/invariance_*.json}. \path{scripts/fairness_report.py} joins \code{nba.cards} to \code{gold.c360\_restricted} on a compliance
machine and caches the rates in \path{data/exports/reports/fairness_data.json}; the separated-customer breakdown in §4 is a
one-off query on the same tables.}{\code{uv run python scripts/invariance\_test.py protected} and \code{all}; \code{uv run python scripts/fairness\_report.py} (with \code{kept serve} stopped), then
\code{uv run python scripts/reportlib.py fairness\_report}}

\end{document}
